\chapter{Lab Solutions and Acceptance Guidance}
\index{lab solutions}\index{acceptance criteria}\index{hands-on project!solutions}%
This appendix provides reference solutions for every chapter's Thursday hands-on project
(Chapters~0--24) and selected end-of-chapter exercises. Each entry gives the approach, the key
commands or code fragments, the evidence that satisfies the chapter's acceptance criteria, and
the failure modes that actually occur. Use it to unblock yourself after a genuine attempt---not
instead of one: the labs are the curriculum, and the AWS Certified Data Engineer~-- Associate
exam measures exactly the judgment they build \cite{awsCertDataEngineer}.

\begin{notebox}
A solution counts only when the chapter's acceptance criteria are met with evidence: row counts,
metric screenshots, ledger rows, drill outputs. ``It ran once on my machine'' is the starting
point of each lab, never the end. Replace every placeholder---bucket names, account IDs, Regions,
role ARNs---with your own before running anything.
\end{notebox}

\section{Chapter 0: Environment Verification}
\subsection{Project: Account, Profile, Guardrails, and the S3 Smoke Test}
\textbf{Approach.} Secure root (MFA), create the Identity Center admin, configure the
\texttt{sandbox} CLI profile, create a budget with alerts, then verify:
\begin{lstlisting}[language=bash,numbers=none]
aws sts get-caller-identity --profile sandbox
aws s3 mb s3://<unique-bucket> ; aws s3 ls ; aws s3 rb s3://<unique-bucket>
\end{lstlisting}
\textbf{Critical evidence.} Caller identity shows the admin (never root); the budget exists with
three thresholds; bucket create/list/delete succeed in the intended Region.
\textbf{Common failure.} SSO token expired mid-lab---re-run \texttt{aws sso login}; or working in
the console's default Region while the CLI profile points elsewhere.

\section{Part I: Core Data Infrastructure and Storage (Chapters 1--4)}
\subsection{Chapter 1: Bucket, Lifecycle, and Events}
\textbf{Approach.} CLI-first: create the bucket, enable versioning, apply a 30-day transition to
Glacier Flexible Retrieval, create the SQS queue, attach the queue policy, register the S3 event
notification, upload objects, and read the events back.
\textbf{Key fragment.}
\begin{lstlisting}[language=bash,numbers=none]
aws s3api put-bucket-versioning --bucket $B --versioning-configuration Status=Enabled
aws s3api put-bucket-lifecycle-configuration --bucket $B --lifecycle-configuration file://lc.json
aws s3api put-bucket-notification-configuration --bucket $B `
  --notification-configuration file://events.json
aws sqs receive-message --queue-url $Q --max-number-of-messages 10
\end{lstlisting}
\textbf{Critical evidence.} The Python checker confirms versioning, lifecycle rules, and the SQS
binding; an uploaded object produces a readable S3 \texttt{ObjectCreated} event.
\textbf{Common failure.} The SQS queue policy missing the \texttt{Service: s3.amazonaws.com}
principal, so the notification configuration is rejected; or bucket names colliding with the
global namespace---append initials plus the account ID.

\subsection{Chapter 2: Aurora Serverless v2 and Regional Failover}
\textbf{Approach.} Create an Aurora PostgreSQL Serverless v2 cluster with a small ACU range,
attach it to an Aurora Global Database with a secondary Region, then rehearse the failover:
switchover, verify the writer moved, switch back.
\textbf{Key fragment.}
\begin{lstlisting}[language=bash,numbers=none]
aws rds create-db-cluster --engine aurora-postgresql --engine-mode provisioned `
  --serverless-v2-scaling-configuration MinCapacity=0.5,MaxCapacity=4 ...
aws rds switchover-global-cluster --global-cluster-identifier $G `
  --target-db-cluster-identifier $SecondaryArn
\end{lstlisting}
\textbf{Critical evidence.} The secondary cluster becomes the writer; a write/read round-trip
succeeds against the new endpoint; the cluster and secrets are deleted afterward.
\textbf{Common failure.} Forgetting that global-cluster teardown is ordered (instances, then
clusters, then the global cluster), or leaving the secondary Region's cluster running and
billing.

\subsection{Chapter 3: Forum Single-Table Schema and GSI}
\textbf{Approach.} Design access patterns first, then create one table holding users, forums,
threads, and posts with generic \texttt{PK}/\texttt{SK} attributes, plus a GSI for user-centric
queries. Insert representative items and run every declared access pattern as a query.
\textbf{Key fragment.}
\begin{lstlisting}[language=bash,numbers=none]
aws dynamodb create-table --table-name forum-single-table-demo `
  --attribute-definitions AttributeName=PK,AttributeType=S ... `
  --key-schema AttributeName=PK,KeyType=HASH AttributeName=SK,KeyType=RANGE `
  --global-secondary-indexes file://gsi.json --billing-mode PAY_PER_REQUEST
\end{lstlisting}
\textbf{Critical evidence.} Each access pattern answered by \texttt{Query} (not \texttt{Scan});
the GSI returns user-centric ordering; item sizes stay small.
\textbf{Common failure.} Modeling separate tables per entity (relational habits) and then trying
to join client-side; or choosing a low-cardinality partition key that concentrates hot traffic.

\subsection{Chapter 4: Redis Cache-Aside for RDS Queries}
\textbf{Approach.} Deploy a small ElastiCache Redis replication group in the sandbox VPC, then a
Python cache-aside layer: check cache, on miss query RDS and store with a TTL. Measure latency
before/after and deliberately invalidate.
\textbf{Key fragment.}
\begin{lstlisting}[language=Python,numbers=none]
cached = r.get(key)
if cached is None:
    rows = cur.execute(sql).fetchall()
    r.setex(key, 300, json.dumps(rows))   # TTL bounds staleness
\end{lstlisting}
\textbf{Critical evidence.} Second identical query served from cache with sub-millisecond reads;
after a source write, the invalidation path (or TTL expiry) makes the change visible.
\textbf{Common failure.} Security-group mismatch (Lambda/app SG not allowed to the cache SG on
6379); or caching without TTL so staleness is unbounded.

\section{Part II: Data Warehousing and Analytics (Chapters 5--8)}
\subsection{Chapter 5: Load a Star Schema into Redshift Serverless}
\textbf{Approach.} Create namespace + workgroup with an attached least-privilege role, generate
ten million synthetic order-line rows as ten JSON files (so COPY parallelizes across slices),
declare the grain, create \texttt{fact\_sales} plus three conformed dimensions, COPY in, then
practice the SCD Type 2 MERGE against a staged customer extract.
\textbf{Critical evidence.} \texttt{SELECT COUNT(*)} on the fact equals the generated row count;
dimension lookups join cleanly; the SCD2 MERGE run twice changes nothing on the second run.
\textbf{Common failure.} The role trust policy missing \texttt{redshift.amazonaws.com}, so COPY
fails authorization; or generating one giant file so COPY serializes on the leader's file
assignment instead of parallelizing.

\subsection{Chapter 6: Benchmarking Distribution Styles}
\textbf{Approach.} Create three copies of the fact table---\texttt{DISTSTYLE KEY} on the join
column, \texttt{EVEN}, and \texttt{ALL}---then run one join-heavy aggregation against each,
record elapsed time, and read \texttt{EXPLAIN} for \texttt{DS\_DIST\_OUTER}/broadcast steps.
\textbf{Critical evidence.} The completed timing table; the KEY plan showing a collocated join;
the ALL plan's storage and load-time cost quantified.
\textbf{Common failure.} Benchmarking on data too small for redistribution to matter
(regenerate a larger dimension); or comparing cold versus cached runs---clear the result cache
or vary the query text.

\subsection{Chapter 7: Crawl, Query, and Federate the Same Data}
\textbf{Approach.} Run a Glue crawler over raw CSVs, aggregate in Athena recording bytes
scanned, CTAS to partitioned Parquet, catalog the result, rerun (bytes scanned collapses), then
create a Spectrum external schema and join lake to warehouse.
\textbf{Key fragment.}
\begin{lstlisting}[language=SQL,numbers=none]
CREATE TABLE sales_parquet WITH (format='PARQUET',
  external_location='s3://.../curated/sales/', partitioned_by=ARRAY['dt'])
AS SELECT *, CAST(order_date AS date) AS dt FROM raw_sales_csv;
\end{lstlisting}
\textbf{Critical evidence.} The bytes-scanned before/after pair (expect one to two orders of
magnitude); the Spectrum join returning correct counts.
\textbf{Common failure.} Crawling a prefix that mixes datasets so the crawler creates spurious
tables; or forgetting the Athena workgroup results bucket.

\subsection{Chapter 8: Hiding the PII Column from Marketing}
\textbf{Approach.} Register the S3 location with Lake Formation, designate the admin, revoke
\texttt{IAMAllowedPrincipals} on database and table, grant the marketing role \texttt{SELECT} on
a column list excluding \texttt{pii\_email}, then verify from both roles in Athena.
\textbf{Critical evidence.} Marketing query on allowed columns succeeds; the same query with
\texttt{pii\_email} fails with a Lake Formation error; the analytics role still sees everything.
\textbf{Common failure.} Testing as the lake admin---every grant passes trivially; or revoking
\texttt{IAMAllowedPrincipals} on the database but not the table, so the filter never binds.

\section{Part III: Batch Processing and ETL Orchestration (Chapters 9--12)}
\subsection{Chapter 9: A PySpark Aggregation on EMR}
\textbf{Approach.} Launch EMR on EC2 with instance fleets (On-Demand core, Spot task), submit a
PySpark step aggregating the Chapter~7 curated Parquet, write results back to S3, inspect the
Spark UI, terminate, and prove the data survived.
\textbf{Critical evidence.} Output Parquet in S3 after cluster termination; Spark UI shows the
distributed aggregation (stages, tasks across executors); Spot interruptions rescheduled tasks
without job failure.
\textbf{Common failure.} Putting durable output on cluster HDFS (lost at termination); or Spot on
core nodes risking HDFS loss---only task nodes are Spot-safe.

\subsection{Chapter 10: A Glue Studio Pipeline into Redshift}
\textbf{Approach.} Visual canvas: two catalog sources (\texttt{week7\_lake.sales\_parquet},
\texttt{week7\_lake.customers}), a join, a null filter, a timestamp transform, a Redshift target.
Then harden the generated PySpark: parameterize, enable bookmarks, run twice, and inject one bad
record.
\textbf{Critical evidence.} Second run with bookmarks processes nothing new; the bad record is
visibly handled (dropped or quarantined, per the transform); row counts reconcile between source
and Redshift.
\textbf{Common failure.} Bookmarks enabled but the sink not idempotent---re-runs still duplicate
rows; or forgetting the Redshift connection's VPC/subnet placement so the job cannot reach the
endpoint.

\subsection{Chapter 11: CDC from RDS PostgreSQL to the Lake}
\textbf{Approach.} Enable logical replication (\texttt{rds.logical\_replication=1} in the
parameter group, reboot, confirm \texttt{wal\_level=logical}), create a DMS replication instance,
source and target endpoints, and a full-load-plus-CDC task. Generate source changes, watch change
files land in S3, alarm on lag.
\textbf{Critical evidence.} INSERT/UPDATE/DELETE on the source appear as S3 change files;
\texttt{CDCLatencySource}/\texttt{CDCLatencyTarget} stay near zero under light load and the
CloudWatch alarm exists.
\textbf{Common failure.} Skipping the parameter-group reboot, so the task fails at CDC start; or
WAL retention too low so a paused task cannot resume and forces a reload.

\subsection{Chapter 12: A Production-Shaped MWAA DAG}
\textbf{Approach.} Deploy the daily-reconciliation DAG to the MWAA DAG bucket:
\texttt{S3KeySensor} $\rightarrow$ \texttt{GlueJobOperator} $\rightarrow$ Athena validation
$\rightarrow$ SNS alert, with \texttt{catchup=False}, retries with delay, a 300-second
\texttt{poke\_interval}, and \texttt{trigger\_rule="one\_failed"} on the alert.
\textbf{Critical evidence.} A forced Glue failure fires the SNS alert with the runbook link;
a normal run lands one ledger row; the execution role is least-privilege to exactly the four
services used.
\textbf{Common failure.} The alert task skipped because the default trigger rule skips on
upstream skip---exactly the runs it exists for; or \texttt{catchup=True} with a past
\texttt{start\_date} backfilling against production on deploy.

\section{Part IV: Streaming and Real-Time Analytics (Chapters 13--16)}
\subsection{Chapter 13: An IoT Temperature Stream}
\textbf{Approach.} Shard math first (50 sensors $\times$ 150 B $\approx$ 7.5 KB/s, 50
records/s $\rightarrow$ one shard suffices; two used to make sharding visible). Python producer
with \texttt{put\_records}, partial-failure handling with backoff, per-sensor partition keys for
ordering; a raw-iterator consumer printing records.
\textbf{Critical evidence.} Records grouped per sensor arrive in order per shard;
\texttt{IteratorAgeMilliseconds} rises when the consumer pauses and falls to zero on catch-up.
\textbf{Common failure.} Ignoring \texttt{FailedRecordCount} (individual records fail inside a
200 batch); or using the sensor \emph{type} as the partition key and hot-sharding one key.

\subsection{Chapter 14: Stream to Parquet via Firehose}
\textbf{Approach.} Delivery stream sourcing from the Chapter~13 Kinesis stream, a Lambda that
uppercases string values, JSON$\rightarrow$Parquet conversion against a Glue table, S3
destination with error prefixes and raw backup enabled.
\textbf{Critical evidence.} Curated Parquet files under the date-partitioned prefix; original
JSON in the backup prefix; a malformed record in the error prefix.
\textbf{Common failure.} The Glue table schema not matching the transformed JSON (conversion
fails silently into the error prefix); or Lambda timeouts from per-record work that belongs in
batch.

\subsection{Chapter 15: A Flink SQL Tumbling-Window Counter}
\textbf{Approach.} Source and sink Kinesis streams; Flink SQL tables with a declared watermark;
a 5-minute \texttt{TUMBLE} count; then deliberately send on-time, late-within-watermark, and
late-beyond-watermark events and observe.
\textbf{Key fragment.}
\begin{lstlisting}[language=SQL,numbers=none]
CREATE TABLE clicks (user_id STRING, page STRING, click_ts TIMESTAMP(3),
  WATERMARK FOR click_ts AS click_ts - INTERVAL '1' MINUTE) WITH (...);
SELECT TUMBLE_START(click_ts, INTERVAL '5' MINUTE) AS w, COUNT(*) AS n
FROM clicks GROUP BY TUMBLE(click_ts, INTERVAL '5' MINUTE);
\end{lstlisting}
\textbf{Critical evidence.} Windows fire on the watermark, not the wall clock; the beyond-bound
late event is dropped (or side-output, if configured) and that behavior is documented.
\textbf{Common failure.} Judging lateness by ingestion time---the table must declare event time
plus watermark or every ``late'' test is meaningless.

\subsection{Chapter 16: A Monitored Sales Dashboard}
\textbf{Approach.} QuickSight dataset over \texttt{fact\_sales} joined to its dimensions; KPI +
trend line + working region filter; SPICE import with a refresh schedule; ML anomaly detection on
daily revenue alerting to email/SNS; the same dataset in direct query for comparison.
\textbf{Critical evidence.} The dashboard filtered per region; the anomaly alert fires on an
injected dip; SPICE vs.\ direct-query latency and freshness notes recorded.
\textbf{Common failure.} Publishing before RLS rules are tested with a non-author identity; or
scheduling SPICE refresh slower than the freshness requirement the dashboard promises.

\section{Part V: Machine Learning and MLOps (Chapters 17--20)}
\subsection{Chapter 17: Churn Prediction in Pure SQL}
\textbf{Approach.} \texttt{RedshiftMLRole} trusted by Redshift with S3 staging and SageMaker
permissions; training query with the leaky column deliberately excluded (and documented);
\texttt{CREATE MODEL} with explicit problem type and F1 objective; score the customer base and
persist a risk table.
\textbf{Critical evidence.} Validation metrics show F1 (not accuracy) as the objective; the risk
table exists and a second scoring run is idempotent; the excluded leaky column is named in the
write-up.
\textbf{Common failure.} Training on a feature only knowable after the outcome (target leakage)
producing a ``too good'' model; or forgetting the staging bucket grant so the export step fails.

\subsection{Chapter 18: Wrangler to Serverless Endpoint}
\textbf{Approach.} Governed S3 dataset read through the Lake Formation grant; Data Wrangler flow
(type fixes, null imputation, one-hot) exported as a Processing Job; Autopilot experiment
including XGBoost candidates; best model to a serverless endpoint; test payload.
\textbf{Critical evidence.} The exported flow reruns without edits; the endpoint scales to zero
and answers a payload; the candidate leaderboard was inspected, not just accepted.
\textbf{Common failure.} The Studio execution role bypassing Lake Formation with broad S3
access---the lab's governance point is then untested; or leaving a real-time endpoint running
when the lab called for serverless.

\subsection{Chapter 19: A SageMaker Pipeline, End to End}
\textbf{Approach.} Four-step pipeline---Preprocess $\rightarrow$ Train $\rightarrow$ Evaluate
$\rightarrow$ Register---with parameters for input data and thresholds, a condition step gating
registration on the metric, executed twice: once passing, once with an injected quality drop.
\textbf{Critical evidence.} One run registers, one run halts at the gate; lineage shows the
registered version's data and training inputs; the approved version is attributable in the
Registry.
\textbf{Common failure.} Step caching masking a real change (inputs changed but fingerprints
did not); or the condition comparing the wrong metric name so the gate never fires.

\subsection{Chapter 20: S3-Triggered Textract Extraction}
\textbf{Approach.} S3 event notification on \texttt{raw/documents/} invoking Lambda; sync
\texttt{AnalyzeDocument} with \texttt{FORMS} for single pages (async + SNS for multi-page);
block structure parsed into key--value JSON; results persisted under
\texttt{extracted/documents/} keyed by source etag for idempotency.
\textbf{Critical evidence.} Re-uploading the same file does not re-bill (etag check skips); a
deliberately corrupt upload lands in the visible failure path; extracted JSON matches the form's
fields.
\textbf{Common failure.} Sending a multi-page PDF to the sync API (it is single-page only); or
persisting raw OCR output into the curated zone---it faithfully contains the document's PII.

\section{Part VI: Security, Governance, and Exam Prep (Chapters 21--24)}
\subsection{Chapter 21: Private S3 Access with No Internet Route}
\textbf{Approach.} VPC with a private subnet and deliberately no internet gateway; S3 gateway
endpoint with a restrictive endpoint policy; EC2 reached via Session Manager; copy a file to the
authorized bucket, then prove a non-authorized bucket is blocked.
\textbf{Critical evidence.} The upload succeeds with no public route in the route table; the
endpoint policy denies the other bucket; \texttt{aws s3 ls} against the unauthorized bucket fails
with an explicit deny.
\textbf{Common failure.} An interface endpoint where a gateway endpoint was meant (hourly cost
plus DNS confusion); or Session Manager failing because the instance profile lacks
\texttt{AmazonSSMManagedInstanceCore}.

\subsection{Chapter 22: A Quality Gate That Halts the Pipeline}
\textbf{Approach.} DQDL ruleset evaluated inside the Glue job between transform and sink:
\texttt{ColumnExists} for structure plus \texttt{Completeness "col" >= 0.95} for content. Pass
writes; fail publishes to SNS and raises. Inject a violating batch.
\textbf{Critical evidence.} The violating batch halts the job, publishes nothing downstream, and
pages via SNS; outcomes recorded to CloudWatch for trending; the next clean batch passes.
\textbf{Common failure.} The gate evaluated but not wired---results logged while the write
proceeds anyway (aspirational quality control); or thresholdless rules paging on noise until the
team learns to ignore them.

\subsection{Chapter 23: Macie Discovery and Automated Quarantine}
\textbf{Approach.} Enable Macie, run a one-time classification job over a test prefix, upload a
synthetic CSV with mock card numbers, then wire EventBridge \texttt{Macie Finding}
$\rightarrow$ Lambda $\rightarrow$ move the object to an Object-Locked quarantine bucket and
write an audit record linking finding ID to action.
\textbf{Critical evidence.} The finding appears with the managed identifier match; the file moves
to quarantine immutably; the audit record resolves finding $\rightarrow$ object $\rightarrow$
action.
\textbf{Common failure.} Deleting instead of quarantining (destroys the evidence the finding
refers to); or testing with real PII---synthetic fixtures only, always.

\subsection{Chapter 24: The Repository Is the Platform}
\textbf{Approach.} One GitHub repository: \texttt{terraform/} (module + Dev/Prod environments
over an S3+DynamoDB backend), \texttt{dbt/} (three models---staging, intermediate, mart---and
five tests against Redshift), and \texttt{.github/workflows/} with OIDC-to-IAM federation.
\textbf{Critical evidence.} A failing dbt test blocks merge; merge deploys to Dev; Prod waits on
named approval; \texttt{terraform plan} on a second machine against the remote backend shows no
drift.
\textbf{Common failure.} Long-lived AWS keys stored in CI secrets instead of OIDC; or state
committed locally so two engineers diverge the same environment.

\section{Selected End-of-Chapter Exercises}
\index{exercises!solutions}%
\subsection{Storage and Warehousing}
\begin{itemize}
    \item \textbf{(Ch.~1, design) Lifecycle for mixed-size objects.} Compute per-object
    transition overhead against the actual size histogram first; exclude objects below 128~KB
    from Glacier transitions and let them expire in Standard-IA instead. The correct answer is
    usually a \emph{filter}, not a fancier rule.
    \item \textbf{(Ch.~5, design) Choose fact types for an order pipeline.} Order \emph{line} is
    the transaction fact; daily inventory is the periodic snapshot; the order itself
    (placed $\rightarrow$ shipped $\rightarrow$ delivered) is the accumulating snapshot. Grains
    differ; never merge them.
    \item \textbf{(Ch.~6, tuning) Diagnose a slow dashboard without touching SQL first.} Check
    WLM queue wait time: if waits dominate, the fix is capacity or concurrency scaling, not
    query rewriting. Only when execution time dominates do sort keys, MVs, and distribution
    enter.
\end{itemize}
\subsection{ETL and Streaming}
\begin{itemize}
    \item \textbf{(Ch.~10, debugging) A bookmarked job still duplicates rows.} Bookmarks only
    deduplicate the \emph{source} side. The sink must be idempotent: overwrite the target
    partition or MERGE by business key.
    \item \textbf{(Ch.~13, design) Replay requirement chooses the service.} Multiple consumers
    and replay rule out SQS FIFO; one consumer, no replay rules out Kinesis. State both
    requirements before naming a service.
    \item \textbf{(Ch.~15, theory) Exactly-once to S3.} Checkpointing plus a two-phase-commit
    (or idempotent-by-key) sink; checkpointing alone only makes \emph{internal} state
    consistent.
\end{itemize}
\subsection{ML, Security, and DataOps}
\begin{itemize}
    \item \textbf{(Ch.~18, design) Endpoint choice for sporadic traffic.} Serverless endpoints:
    scale to zero, pay per request, accept cold starts. Real-time only for steady SLA-bound
    traffic; Batch Transform for bulk offline scoring.
    \item \textbf{(Ch.~21, design) Close the exfiltration path.} Security groups alone cannot
    stop writes to a stranger's bucket: endpoint policies restrict destinations, SCPs cap the
    account, KMS key policies make copied data unreadable. All three layers.
    \item \textbf{(Ch.~24, design) Exam-prep plan for a teammate blank on streaming and
    security.} Rebuild, do not re-read: the Kinesis shard-math and Flink watermark labs
    (Chapters~13, 15), then the Lake Formation grant and VPC endpoint labs (Chapters~8, 21),
    each followed by the chapter exercises attempted cold.
\end{itemize}

\begin{tipbox}
When your solution diverges from this appendix, the tie-breaker is the acceptance criteria in the
chapter, not the reference here: labs admit many correct implementations, but only one definition
of done.
\end{tipbox}
